\documentclass[11pt]{article}
\usepackage[margin=2.5cm]{geometry}
\usepackage{times}
\usepackage{longtable}
\usepackage{array}
\usepackage{hyperref}
\usepackage{listings}
\usepackage{xcolor}
\usepackage{colortbl}
\usepackage{parskip}
\usepackage{titlesec}
\usepackage{enumitem}
\usepackage{fancyhdr}
\usepackage{tocloft}
\setlength{\cftbeforesecskip}{2pt}
\setlength{\cftbeforesubsecskip}{0pt}
\setlength{\cftbeforesubsubsecskip}{0pt}
\setlength{\headheight}{15pt}
\setlength{\tabcolsep}{6pt}
\renewcommand{\arraystretch}{1.3}
\setcounter{secnumdepth}{0}
\setlength{\LTleft}{0pt}
\setlength{\LTright}{0pt}

% ─── Hyperlink colours ────────────────────────────────────────────────────────
\hypersetup{
  colorlinks=true,
  linkcolor=blue,
  urlcolor=blue
}

% ─── Colours ──────────────────────────────────────────────────────────────────
\definecolor{tabheader}{RGB}{189,215,238}
\definecolor{codebg}{RGB}{189,215,238}
\definecolor{keywordcolor}{RGB}{0,100,200}
\definecolor{stringcolor}{RGB}{180,50,0}
\definecolor{numbercolor}{RGB}{0,128,0}

% ─── Code listing style ───────────────────────────────────────────────────────
\lstdefinelanguage{json}{
  basicstyle=\ttfamily\small,
  backgroundcolor=\color{codebg},
  frame=none,
  showstringspaces=false,
  breaklines=true,
  postbreak=\mbox{\textcolor{gray}{$\hookrightarrow$}\space},
  literate=
    *{0}{{{\color{numbercolor}0}}}{1}
     {1}{{{\color{numbercolor}1}}}{1}
     {2}{{{\color{numbercolor}2}}}{1}
     {3}{{{\color{numbercolor}3}}}{1}
     {4}{{{\color{numbercolor}4}}}{1}
     {5}{{{\color{numbercolor}5}}}{1}
     {6}{{{\color{numbercolor}6}}}{1}
     {7}{{{\color{numbercolor}7}}}{1}
     {8}{{{\color{numbercolor}8}}}{1}
     {9}{{{\color{numbercolor}9}}}{1}
     {:}{{{\color{black}{:}}}}{1}
     {,}{{{\color{black}{,}}}}{1}
     {\{}{{{\color{black}{\{}}}}{1}
     {\}}{{{\color{black}{\}}}}}{1}
     {[}{{{\color{black}{[}}}}{1}
     {]}{{{\color{black}{]}}}}{1}
     {"true"}{{{\color{numbercolor}"true"}}}{6}
     {"false"}{{{\color{numbercolor}"false"}}}{7}
     {"null"}{{{\color{numbercolor}"null"}}}{6}
}

% ─── Section title formatting ─────────────────────────────────────────────────
\titleformat{\section}{\large\bfseries}{}{0em}{}
\titleformat{\subsection}{\normalsize\bfseries}{}{0em}{}
\titleformat{\subsubsection}{\normalsize\bfseries}{}{0em}{}

% ─── Table column types ───────────────────────────────────────────────────────
\newcolumntype{L}[1]{>{\raggedright\arraybackslash}p{#1}}
\newcolumntype{H}[1]{>{\columncolor{tabheader}\raggedright\arraybackslash}p{#1}}

% ─── Header / Footer ─────────────────────────────────────────────────────────
\pagestyle{fancy}
\fancyhf{}
\renewcommand{\headrulewidth}{0pt}
\fancyfoot[C]{\thepage}

% ─── Document ─────────────────────────────────────────────────────────────────
\begin{document}

% ─── Cover Page ───────────────────────────────────────────────────────────────
\thispagestyle{empty}
\begin{center}
  \vspace*{2cm}
  {\LARGE\bfseries TECHNICAL DOCUMENTATION}\\[0.8em]
  {\Large\bfseries AI Corrector}\\[0.6em]
  {\Large\bfseries Version 1.0.0}\\[2em]
  July 2, 2026
\end{center}

% ─── Table of Contents ────────────────────────────────────────────────────────
\hypersetup{linkcolor=black}
{\setlength{\parskip}{0pt}\tableofcontents}
\hypersetup{linkcolor=blue}
\newpage

% ══════════════════════════════════════════════════════════════════════════════
\section{Changelog}
% ══════════════════════════════════════════════════════════════════════════════

\begin{itemize}[leftmargin=2em]
  \item \textbf{02-07-2026}

    \textbf{Steven Lie}

    Created initial documentation.
\end{itemize}

% ══════════════════════════════════════════════════════════════════════════════
\section{Base URL}
% ══════════════════════════════════════════════════════════════════════════════

\begin{tabular}{|H{3cm}|p{12.1cm}|}
  \hline
  Development & \texttt{http://localhost:8000} \\
  \hline
  UAT         & \texttt{https://\textit{<uat-host>}:8000} \\
  \hline
  Production  & \texttt{https://\textit{<prod-host>}:8000} \\
  \hline
\end{tabular}

% ══════════════════════════════════════════════════════════════════════════════
\section{Authentication}
% ══════════════════════════════════════════════════════════════════════════════

When the \texttt{API\_KEY} environment variable is set, all endpoints (except
\texttt{/}, \texttt{/health}, \texttt{/docs}, \texttt{/redoc}, and
\texttt{/openapi.json}) require the following request header:

\begin{tabular}{|H{4cm}|p{11.1cm}|}
  \hline
  Header Name & \texttt{X-API-Key} \\
  \hline
  Value       & Your configured API key \\
  \hline
\end{tabular}

% ══════════════════════════════════════════════════════════════════════════════
\section{Health Check}
% ══════════════════════════════════════════════════════════════════════════════

\begin{tabular}{|H{3cm}|p{12.1cm}|}
  \hline
  URL           & \texttt{\{\{BaseURL\}\}/health} \\
  \hline
  Method        & GET \\
  \hline
  Authorization & None \\
  \hline
\end{tabular}

\subsection{Response Fields}

\begin{longtable}{|p{3.5cm}|p{2cm}|p{9.2cm}|}
  \hline
  \rowcolor{tabheader}
  \textbf{Field Name} & \textbf{Type} & \textbf{Note} \\
  \hline
  \endhead
  status  & string & Overall service status: \texttt{"ok"} or \texttt{"degraded"} \\
  \hline
  checks  & object & See \hyperref[sec:healthchecks]{Health Checks} \\
  \hline
\end{longtable}

\subsubsection{Health Checks}
\label{sec:healthchecks}

\begin{longtable}{|p{4.5cm}|p{2cm}|p{8.2cm}|}
  \hline
  \rowcolor{tabheader}
  \textbf{Field Name} & \textbf{Type} & \textbf{Note} \\
  \hline
  \endhead
  vectordb                 & object  & Azure AI Search connectivity status \\
  \hline
  vectordb.status          & string  & \texttt{"ok"} or \texttt{"error"} \\
  \hline
  vectordb.document\_count & integer & Number of indexed documents \\
  \hline
  vectordb.error           & string  & Error message if status is \texttt{"error"} \\
  \hline
\end{longtable}

\subsection{Example Response}

\begin{lstlisting}[language=json]
{
    "status": "ok",
    "checks": {
        "vectordb": {
            "status": "ok",
            "document_count": 150
        }
    }
}
\end{lstlisting}

% ══════════════════════════════════════════════════════════════════════════════
\section{Feed — Course Material Ingestion}
% ══════════════════════════════════════════════════════════════════════════════

% ─── POST /feed ───────────────────────────────────────────────────────────────
\subsection{Upload Files}

\begin{tabular}{|H{3cm}|p{12.1cm}|}
  \hline
  URL           & \texttt{\{\{BaseURL\}\}/feed} \\
  \hline
  Method        & POST \\
  \hline
  Content-Type  & \texttt{multipart/form-data} \\
  \hline
  Authorization & X-API-Key \\
  \hline
\end{tabular}

Upload one or more course material files (PDF, PPT, PPTX, DOCX, TXT). Files
are parsed, chunked, embedded, and stored in Azure AI Search.

\subsubsection{Request Fields}

\begin{longtable}{|p{3.5cm}|p{2cm}|p{2cm}|p{6.8cm}|}
  \hline
  \rowcolor{tabheader}
  \textbf{Field Name} & \textbf{Type} & \textbf{Required} & \textbf{Note} \\
  \hline
  \endhead
  course\_code & string & Yes & Course code to associate materials with, e.g.\ \texttt{COMP6100} \\
  \hline
  files        & file[] & Yes & One or more files (PDF, PPT, PPTX, DOCX, TXT) \\
  \hline
\end{longtable}

\subsubsection{Response Fields}

\begin{longtable}{|p{5cm}|p{2cm}|p{7.7cm}|}
  \hline
  \rowcolor{tabheader}
  \textbf{Field Name} & \textbf{Type} & \textbf{Note} \\
  \hline
  \endhead
  status             & string  & Always \texttt{"completed"} \\
  \hline
  results            & array   & Per-file result list; see Result Item \\
  \hline
  token\_usage       & object  & Aggregated token usage; see \hyperref[sec:tokenusage]{Token Usage} \\
  \hline
  \multicolumn{3}{|l|}{\textit{Result Item}} \\
  \hline
  results[].status              & string  & \texttt{"success"} or \texttt{"failed"} \\
  \hline
  results[].filename            & string  & Original filename \\
  \hline
  results[].total\_chunks\_saved & integer & Number of chunks stored \\
  \hline
  results[].token\_usage        & object  & Per-file token usage \\
  \hline
  results[].error               & string  & Error message (only when \texttt{"failed"}) \\
  \hline
\end{longtable}

\subsubsection{Example Response}

\begin{lstlisting}[language=json]
{
    "status": "completed",
    "results": [
        {
            "status": "success",
            "filename": "lecture1.pdf",
            "total_chunks_saved": 45,
            "token_usage": {
                "embedding_tokens": 2000,
                "embedding_cost_usd": 0.044,
                "vision_tokens": 500,
                "vision_cost_usd": 1.25,
                "total_cost_usd": 1.294
            }
        }
    ],
    "token_usage": {
        "embedding_tokens": 2000,
        "embedding_cost_usd": 0.044,
        "vision_tokens": 500,
        "vision_cost_usd": 1.25,
        "total_cost_usd": 1.294
    }
}
\end{lstlisting}

% ─── POST /feed-url ───────────────────────────────────────────────────────────
\subsection{Upload from Single URL}

\begin{tabular}{|H{3cm}|p{12.1cm}|}
  \hline
  URL           & \texttt{\{\{BaseURL\}\}/feed-url} \\
  \hline
  Method        & POST \\
  \hline
  Content-Type  & \texttt{application/json} \\
  \hline
  Authorization & X-API-Key \\
  \hline
\end{tabular}

Download a course material file from a URL and ingest it into the vector
database.

\subsubsection{Request Fields}

\begin{longtable}{|p{3.5cm}|p{2cm}|p{2cm}|p{6.8cm}|}
  \hline
  \rowcolor{tabheader}
  \textbf{Field Name} & \textbf{Type} & \textbf{Required} & \textbf{Note} \\
  \hline
  \endhead
  url          & string & Yes & Publicly accessible file URL \\
  \hline
  course\_code & string & Yes & Course code to associate material with \\
  \hline
  token        & string & No  & Bearer token for protected URLs \\
  \hline
\end{longtable}

\subsubsection{Example Request}

\begin{lstlisting}[language=json]
{
    "url": "https://example.com/lecture1.pdf",
    "course_code": "COMP6100",
    "token": ""
}
\end{lstlisting}

\subsubsection{Response Fields}

\begin{longtable}{|p{4cm}|p{2cm}|p{8.7cm}|}
  \hline
  \rowcolor{tabheader}
  \textbf{Field Name} & \textbf{Type} & \textbf{Note} \\
  \hline
  \endhead
  status               & string  & \texttt{"success"} or \texttt{"failed"} \\
  \hline
  message              & string  & Human-readable result message \\
  \hline
  total\_chunks\_saved & integer & Number of chunks stored \\
  \hline
  token\_usage         & object  & See \hyperref[sec:tokenusage]{Token Usage} \\
  \hline
\end{longtable}

\subsubsection{Example Response}

\begin{lstlisting}[language=json]
{
    "status": "success",
    "message": "'lecture1.pdf' inserted from URL",
    "total_chunks_saved": 45,
    "token_usage": {
        "embedding_tokens": 2000,
        "embedding_cost_usd": 0.044,
        "vision_tokens": 500,
        "vision_cost_usd": 1.25,
        "total_cost_usd": 1.294
    }
}
\end{lstlisting}

% ─── POST /feed-urls ──────────────────────────────────────────────────────────
\subsection{Upload from Multiple URLs}

\begin{tabular}{|H{3cm}|p{12.1cm}|}
  \hline
  URL           & \texttt{\{\{BaseURL\}\}/feed-urls} \\
  \hline
  Method        & POST \\
  \hline
  Content-Type  & \texttt{application/json} \\
  \hline
  Authorization & X-API-Key \\
  \hline
\end{tabular}

Download and ingest multiple files concurrently. Individual failures do not
fail the entire request.

\subsubsection{Request Fields}

\begin{longtable}{|p{3.5cm}|p{2cm}|p{2cm}|p{6.8cm}|}
  \hline
  \rowcolor{tabheader}
  \textbf{Field Name} & \textbf{Type} & \textbf{Required} & \textbf{Note} \\
  \hline
  \endhead
  urls         & string[] & Yes & List of file URLs to ingest \\
  \hline
  course\_code & string   & Yes & Course code to associate materials with \\
  \hline
  token        & string   & No  & Bearer token for protected URLs \\
  \hline
\end{longtable}

\subsubsection{Example Request}

\begin{lstlisting}[language=json]
{
    "urls": [
        "https://example.com/lecture1.pdf",
        "https://example.com/lecture2.pptx"
    ],
    "course_code": "COMP6100",
    "token": ""
}
\end{lstlisting}

\subsubsection{Response Fields}

\begin{longtable}{|p{4cm}|p{2cm}|p{8.7cm}|}
  \hline
  \rowcolor{tabheader}
  \textbf{Field Name} & \textbf{Type} & \textbf{Note} \\
  \hline
  \endhead
  status       & string & Always \texttt{"completed"} \\
  \hline
  results      & array  & Per-file result list (same schema as \texttt{POST /feed}) \\
  \hline
  token\_usage & object & Aggregated token usage; see \hyperref[sec:tokenusage]{Token Usage} \\
  \hline
\end{longtable}

% ══════════════════════════════════════════════════════════════════════════════
\section{Assessment — Student Answer Evaluation}
% ══════════════════════════════════════════════════════════════════════════════

% ─── POST /assess ─────────────────────────────────────────────────────────────
\subsection{Evaluate Single Student}

\begin{tabular}{|H{3cm}|p{12.1cm}|}
  \hline
  URL           & \texttt{\{\{BaseURL\}\}/assess} \\
  \hline
  Method        & POST \\
  \hline
  Content-Type  & \texttt{multipart/form-data} \\
  \hline
  Authorization & X-API-Key \\
  \hline
\end{tabular}

Evaluate one student's answer against a question and rubric. When
\texttt{use\_key\_answer} is \texttt{true}, the provided key answer is used as
reference; otherwise the system retrieves relevant context from the vector
database.

\subsubsection{Request Fields}

\begin{longtable}{|p{3.8cm}|p{2cm}|p{1.8cm}|p{6.7cm}|}
  \hline
  \rowcolor{tabheader}
  \textbf{Field Name} & \textbf{Type} & \textbf{Required} & \textbf{Note} \\
  \hline
  \endhead
  question               & string  & Yes & The exam question \\
  \hline
  student\_answer        & string  & Yes & Student's answer text, or URL to document (Google Docs, PDF, web article) \\
  \hline
  student\_answer\_token & string  & No  & Bearer token for protected answer URLs \\
  \hline
  rubric                 & string  & No  & Rubric as a JSON array string; see \hyperref[sec:rubricitem]{Rubric Item} \\
  \hline
  course\_code           & string  & No  & Course code for vector DB context (used when \texttt{use\_key\_answer=false}) \\
  \hline
  use\_key\_answer       & boolean & No  & \texttt{true}: use \texttt{key\_answer\_text} / file as reference. \texttt{false}: retrieve from vector DB. Default: \texttt{true} \\
  \hline
  key\_answer\_text      & string  & No  & Reference answer text \\
  \hline
  key\_answer\_file      & file    & No  & Reference answer file (PDF, PPT, PPTX, DOCX, TXT) \\
  \hline
\end{longtable}

\subsubsection{Example Request}

\begin{lstlisting}[language=json]
{
    "question": "Jelaskan konsep machine learning!",
    "student_answer": "Machine learning adalah subset dari AI yang ...",
    "rubric": "[{\"minScore\":0,\"maxScore\":64,\"proficiency\":\"Poor\",\"criteria\":\"Unable to explain the concept\"},{\"minScore\":85,\"maxScore\":100,\"proficiency\":\"Excellent\",\"criteria\":\"Explains concept with clear examples\"}]",
    "use_key_answer": true,
    "key_answer_text": "Machine learning adalah cabang dari kecerdasan buatan ..."
}
\end{lstlisting}

\subsubsection{Response Fields}

\begin{longtable}{|p{4cm}|p{2cm}|p{8.7cm}|}
  \hline
  \rowcolor{tabheader}
  \textbf{Field Name} & \textbf{Type} & \textbf{Note} \\
  \hline
  \endhead
  status             & string & \texttt{"success"} or \texttt{"error"} \\
  \hline
  student\_answer    & string & Resolved plain text of the student's answer \\
  \hline
  retrieved\_sources & array  & Vector DB chunks used as context; see \hyperref[sec:retrievedsource]{Retrieved Source} \\
  \hline
  evaluation         & object & AI evaluation result; see \hyperref[sec:evaluationresult]{Evaluation Result} \\
  \hline
  token\_usage       & object & See \hyperref[sec:tokenusage]{Token Usage} \\
  \hline
\end{longtable}

\subsubsection{Example Response}

\begin{lstlisting}[language=json]
{
    "status": "success",
    "student_answer": "Machine learning adalah subset dari AI ...",
    "retrieved_sources": [],
    "evaluation": {
        "reasoning": "Jawaban mencakup definisi dan konsep dasar dengan baik.",
        "score": 85.5,
        "confidence": 92,
        "feedback": "Tambahkan contoh aplikasi praktis untuk skor lebih tinggi.",
        "sources": []
    },
    "token_usage": {
        "embedding_tokens": 1500,
        "embedding_cost_usd": 0.033,
        "vision_tokens": 0,
        "vision_cost_usd": 0.0,
        "completion_input_tokens": 2000,
        "completion_output_tokens": 500,
        "completion_cost_usd": 3.25,
        "total_cost_usd": 3.283
    }
}
\end{lstlisting}

% ─── POST /assess-batch ───────────────────────────────────────────────────────
\subsection{Evaluate Batch — Single Question}

\begin{tabular}{|H{3cm}|p{12.1cm}|}
  \hline
  URL           & \texttt{\{\{BaseURL\}\}/assess-batch} \\
  \hline
  Method        & POST \\
  \hline
  Content-Type  & \texttt{application/json} \\
  \hline
  Authorization & X-API-Key \\
  \hline
\end{tabular}

Evaluate multiple students for a single question in parallel. Individual
student failures do not fail the entire request.

\subsubsection{Request Fields}

\begin{longtable}{|p{3.5cm}|p{2cm}|p{2cm}|p{6.8cm}|}
  \hline
  \rowcolor{tabheader}
  \textbf{Field Name} & \textbf{Type} & \textbf{Required} & \textbf{Note} \\
  \hline
  \endhead
  question         & string  & Yes & The exam question \\
  \hline
  rubric           & array   & No  & Array of Rubric Item objects; see \hyperref[sec:rubricitem]{Rubric Item} \\
  \hline
  course\_code     & string  & No  & Course code for vector DB context \\
  \hline
  use\_key\_answer & boolean & No  & Use key answer or vector DB context. Default: \texttt{false} \\
  \hline
  key\_answer      & string  & No  & Reference answer text \\
  \hline
  students         & array   & Yes & Array of Student Answer objects; see \hyperref[sec:studentanswer]{Student Answer} \\
  \hline
\end{longtable}

\subsubsection{Example Request}

\begin{lstlisting}[language=json]
{
    "question": "Jelaskan apa yang dimaksud dengan machine learning!",
    "rubric": [
        {
            "minScore": 0,
            "maxScore": 64,
            "proficiency": "Poor",
            "criteria": "Unable to explain the concept"
        },
        {
            "minScore": 85,
            "maxScore": 100,
            "proficiency": "Excellent",
            "criteria": "Explains concept with clear and practical examples"
        }
    ],
    "course_code": "COMP6100",
    "use_key_answer": false,
    "key_answer": "",
    "students": [
        {
            "student_id": "2501001234",
            "answer": "Machine learning adalah subset dari AI ...",
            "token": null
        },
        {
            "student_id": "2501001235",
            "answer": "https://docs.google.com/document/d/ABC123/edit",
            "token": null
        }
    ]
}
\end{lstlisting}

\subsubsection{Response Fields}

\begin{longtable}{|p{5cm}|p{2cm}|p{7.7cm}|}
  \hline
  \rowcolor{tabheader}
  \textbf{Field Name} & \textbf{Type} & \textbf{Note} \\
  \hline
  \endhead
  status                    & string & \texttt{"success"} \\
  \hline
  retrieved\_sources        & array  & Shared vector DB context chunks; see \hyperref[sec:retrievedsource]{Retrieved Source} \\
  \hline
  results                   & array  & Per-student result list \\
  \hline
  results[].student\_id     & string & Student identifier \\
  \hline
  results[].status          & string & \texttt{"success"} or \texttt{"error"} \\
  \hline
  results[].evaluation      & object & AI evaluation; see \hyperref[sec:evaluationresult]{Evaluation Result} (null on error) \\
  \hline
  results[].student\_answer & string & Resolved student answer text (null on error) \\
  \hline
  results[].error           & string & Error message (only when \texttt{"error"}) \\
  \hline
  token\_usage              & object & Aggregated token usage; see \hyperref[sec:tokenusage]{Token Usage} \\
  \hline
\end{longtable}

% ─── POST /assess-batch-multi ─────────────────────────────────────────────────
\subsection{Evaluate Batch — Multiple Questions}

\begin{tabular}{|H{3cm}|p{12.1cm}|}
  \hline
  URL           & \texttt{\{\{BaseURL\}\}/assess-batch-multi} \\
  \hline
  Method        & POST \\
  \hline
  Content-Type  & \texttt{application/json} \\
  \hline
  Authorization & X-API-Key \\
  \hline
\end{tabular}

Evaluate multiple questions concurrently, each with its own set of students,
rubric, and context source. The request body is an array of Batch Assess
Request objects.

\subsubsection{Request Fields}

The request body is a JSON \textbf{array}. Each element is a
\texttt{BatchAssessRequest} with the same fields as
\texttt{POST /assess-batch}.

\subsubsection{Example Request}

\begin{lstlisting}[language=json]
[
    {
        "question": "Jelaskan machine learning!",
        "rubric": [],
        "course_code": "COMP6100",
        "use_key_answer": false,
        "key_answer": "",
        "students": [
            { "student_id": "2501001234", "answer": "...", "token": null }
        ]
    },
    {
        "question": "Jelaskan deep learning!",
        "rubric": [],
        "course_code": "COMP6100",
        "use_key_answer": true,
        "key_answer": "Deep learning adalah subset dari machine learning ...",
        "students": [
            { "student_id": "2501001234", "answer": "...", "token": null }
        ]
    }
]
\end{lstlisting}

\subsubsection{Response Fields}

\begin{longtable}{|p{5cm}|p{2cm}|p{7.7cm}|}
  \hline
  \rowcolor{tabheader}
  \textbf{Field Name} & \textbf{Type} & \textbf{Note} \\
  \hline
  \endhead
  status                       & string & \texttt{"success"} \\
  \hline
  results                      & array  & Per-question result list \\
  \hline
  results[].question           & string & The original question text \\
  \hline
  results[].retrieved\_sources & array  & Vector DB context for this question \\
  \hline
  results[].results            & array  & Per-student results (same schema as \texttt{/assess-batch}) \\
  \hline
  token\_usage                 & object & Aggregated token usage across all questions; see \hyperref[sec:tokenusage]{Token Usage} \\
  \hline
\end{longtable}

% ══════════════════════════════════════════════════════════════════════════════
\section{Shared Schemas}
% ══════════════════════════════════════════════════════════════════════════════

\subsection{Rubric Item}
\label{sec:rubricitem}

\begin{longtable}{|p{3.5cm}|p{2cm}|p{9.2cm}|}
  \hline
  \rowcolor{tabheader}
  \textbf{Field Name} & \textbf{Type} & \textbf{Note} \\
  \hline
  \endhead
  minScore    & number & Minimum score threshold for this level \\
  \hline
  maxScore    & number & Maximum score threshold for this level \\
  \hline
  proficiency & string & Proficiency label, e.g.\ \texttt{"Poor"}, \texttt{"Good"}, \texttt{"Excellent"} \\
  \hline
  criteria    & string & Human-readable description of what earns this level \\
  \hline
\end{longtable}

\subsection{Student Answer}
\label{sec:studentanswer}

\begin{longtable}{|p{3.5cm}|p{2cm}|p{9.2cm}|}
  \hline
  \rowcolor{tabheader}
  \textbf{Field Name} & \textbf{Type} & \textbf{Note} \\
  \hline
  \endhead
  student\_id & string & Unique student identifier \\
  \hline
  answer      & string & Answer text or URL to a document / Google Docs link \\
  \hline
  token       & string & Bearer token for protected answer URLs (optional) \\
  \hline
\end{longtable}

\subsection{Evaluation Result}
\label{sec:evaluationresult}

\begin{longtable}{|p{3.5cm}|p{2cm}|p{9.2cm}|}
  \hline
  \rowcolor{tabheader}
  \textbf{Field Name} & \textbf{Type} & \textbf{Note} \\
  \hline
  \endhead
  reasoning  & string  & AI's step-by-step scoring rationale \\
  \hline
  score      & number  & Numeric score (range determined by rubric) \\
  \hline
  confidence & integer & AI confidence in the score (0--100) \\
  \hline
  feedback   & string  & Constructive feedback for the student \\
  \hline
  sources    & array   & Web sources cited by the AI (if web search was used) \\
  \hline
\end{longtable}

\subsection{Retrieved Source}
\label{sec:retrievedsource}

\begin{longtable}{|p{3.5cm}|p{2cm}|p{9.2cm}|}
  \hline
  \rowcolor{tabheader}
  \textbf{Field Name} & \textbf{Type} & \textbf{Note} \\
  \hline
  \endhead
  source  & string  & Source filename \\
  \hline
  page    & integer & Page or slide number \\
  \hline
  content & string  & Text content of the retrieved chunk \\
  \hline
\end{longtable}

\subsection{Token Usage}
\label{sec:tokenusage}

\begin{longtable}{|p{4.5cm}|p{2cm}|p{8.2cm}|}
  \hline
  \rowcolor{tabheader}
  \textbf{Field Name} & \textbf{Type} & \textbf{Note} \\
  \hline
  \endhead
  embedding\_tokens         & integer & Tokens used for embedding \\
  \hline
  embedding\_cost\_usd      & number  & Embedding API cost (USD) \\
  \hline
  vision\_tokens            & integer & Tokens used by vision model \\
  \hline
  vision\_cost\_usd         & number  & Vision API cost (USD) \\
  \hline
  completion\_input\_tokens  & integer & LLM prompt tokens (assessment only) \\
  \hline
  completion\_output\_tokens & integer & LLM completion tokens (assessment only) \\
  \hline
  completion\_cost\_usd     & number  & LLM API cost (USD) (assessment only) \\
  \hline
  total\_cost\_usd          & number  & Total API cost (USD) \\
  \hline
\end{longtable}

\end{document}
